\documentclass{article}
\usepackage[utf8]{inputenc}
\usepackage{hyperref}
\usepackage{url}
% \usepackage{apalike}
% \documentclass[NETN,manuscript]{stjour-new}
% \articletype{Research}
\usepackage{graphicx}
\graphicspath{ {figs/} }
% \documentclass{article} % This command is used to set the type of document you are working on such as an article, book, or presentation

\usepackage{geometry} % This package allows the editing of the page layout
\usepackage{amsmath}  % This package allows the use of a large range of mathematical formula, commands, and symbols
\usepackage{graphicx}  % This package allows the importing of images

\newcommand{\question}[2][]{
    \begin{flushleft}
    \textbf{Question #1}: \textit{#2}
    \end{flushleft}
}
\newcommand{\sol}{\textbf{Solution}:} %Use if you want a boldface solution line
\newcommand{\maketitletwo}[2][]{
    \begin{center}
        \Large{\textbf{Assignment #1}}

        Course Title

        \vspace{5pt}
        \normalsize{Matthew Frenkel}

        \today
        \vspace{15pt}
    \end{center}
}
\usepackage{geometry}
\usepackage{amsmath}
\usepackage{graphicx}
\usepackage{natbib}
% \usepackage[backend=biber,
%             style=authoryear,
%             natbib=true,
% ]{biblatex}
% \addbibresource{cover.bib}
\usepackage{fancyhdr}
% \usepackage{natbib}
% \usepackage[font=footnotesize,labelfont=bf]{caption}
% \addbibresource{rmt.bib}
% \bibliographystyle{plain}
% \bibstyle{plain}
% \bibliography{cover_letter}

\begin{document}
\bibliographystyle{plain}
\fancypagestyle{nofooter}{\fancyfoot{}}
\thispagestyle{nofooter}  %  https://tex.stackexchange.com/questions/136384/remove-footer-without-removing-the-header
\setlength{\headheight}{30.0pt}
\fancyhead[C]{\includegraphics[height=0.5in]{stfx_text.png}}
\fancyhead[L]{\includegraphics[height=0.5in]{crest.png}}
\fancyhead[R]{
    \sffamily
    \scriptsize
    4130 University Avenue \\
    B2G 2W5 \\
    Antigonish, Nova Scotia \\
    Canada
}
\normalsize

\title{Cover Letter}
% \begin{center}
% \includegraphics[width=4in]{stfx_logo.png}
% \end{center}


% \author[Derek Berger, Gurpreet Matharoo, Jacob Levman]{
%     Derek M. Berger\affil{1},
%     Gurpreet S. Matharoo\affil{2}
%     \and Jacob Levman\affil{3}
% }
% \affiliation{1}{Computer Science, St. Francis Xavier University, Antigonish, Canada, dberger@stfx.ca}
% \affiliation{2}{Physics, St. Francis Xavier University, Antigonish, Canada, gmatharo@stfx.ca}
% \affiliation{3}{Computer Science, St. Francis Xavier University, Antigonish, Canada, jlevman@stfx.ca}


% \keywords{random matrix theory, feature extraction, functional connectivity,
% functional magnetic resonance imaging, voxelwise connectivity, predictive analysis}

\noindent
\today, \\
To the Many Editors of Network Neuroscience, \\

We hope you will consider our manuscript \emph{Random Matrix Theory Tools for the Analysis of
Functional Magnetic Resonance Imaging Examinations} for publication in Network Neuroscience as a
\emph{Research} article. In it, we apply mathematical tools from random matrix theory (RMT) to
analyze a wide variety of functional magnetic resonance imaging (fMRI) data, and also release
open-source software to facilitate such analyses with a wider variety of data beyond fMRI.

Numerous fMRI studies examine correlations amongst regions of interest (ROI) in the brain. These
``functional connectivity'' analyses typically use ROIs or seed regions in order to simplify
connectivity computations \citep{huettelFunctionalMagneticResonance2004}. Instead, we use
statistics from RMT to summarize \emph{all} voxelwise correlations with no reductions. In
addition, rather than simply reporting findings which are statistically significant, we report all
analyses from a cross-validated predictive framework, where multiple machine learning models are
compared against both each other and proper baselines (e.g. naive / dummy classifiers).

We show that, across a variety of data and with fairly small data, and without any model tuning,
that these RMT reductions have real predictive utility and potential. This suggests RMT could
be a powerful and efficient method for summarizing functional connectivity information. As much of the
research in network neuroscience involves fMRI, functional connectivity, and statistical physics, as
per the \href{https://direct.mit.edu/netn/pages/submission-guidelines#aims}{Aims and Scope} of \emph{Network
Neuroscience}, we believe our manuscript to be an excellent fit for submission to your journal.

Should you consider to review our manuscript, we have suggested a number of potential reviewers on
your submission website that we believe would be suitable due to expertise in RMT and similar
interdisciplinary research. This paper is not under consideration at any other journal (although a
poster of some preliminary results of our study was presented at the International Symposium for
Magnetic Resonance in Medicine 2021 annual conference \citep{bergerOpenSourceRandom2021}). Our lab
has REB ethics approval for conducting retrospective analyses (although we are unaware of the need,
this being an analysis of anonymized public data). All authors approve the manuscript in its
submitted form, and have no conflicts of interest, financial or otherwise.



% \begin{itemize}
% \item Sarika Jalan (sarika@iiti.ac.in) \\
% \small
% ~~~ Centre for Biosciences and Biomedical Engineering, IIT Indore India
% \normalsize
% \item Javeria Hashmi (javeria.hashmi@dal.ca) \\
% \small
% ~~~ Department of Anesthesia, Pain Management and Perioperative Medicine Dalhousie University, Canada
% \normalsize
% \item Nivedita Deo (ndeo@physics.du.ac.in) \\
% \small
% ~~~ Department of Physics and Astrophysics University of Delhi, India
% \normalsize
% \item Guillermo A. Cecchi (gcecchi@us.ibm.com) \\
% \small
% ~~~ Thomas J. Watson Research Center, Yorktown Heights, NY USA
% \normalsize
% \item Randy McIntosh (rmcintosh@research.baycrest.org) \\
% \small
% ~~~ Rotman Research Institute at Baycrest Health Sciences, Toronto, Canada
% \normalsize
% \end{itemize}


\bigskip

\noindent
Sincerely, \\

\noindent
\textbf{Corresponding Author}\\
Jacob Levman, PhD (jlevman@stfx.ca) \\
Canada Research Chair in Bioinformatics \\
Department of Computer Science \\
St. Francis Xavier University \\

\noindent
\textbf{Additional Contacts} [should the corresponding author be unavailable]\\
Derek Berger (dberger@stfx.ca), Department of Computer Science, St. Francis Xavier University\\
Gurpreet S. Matharoo (gmatharo@stfx.ca), Department of Physics, St. Francis Xavier University

% \printbibliography % biblatex
\renewcommand{\bibsection}{}  % https://tex.stackexchange.com/questions/132646/how-to-remove-the-references-title
\bibliography{cover.bib}  % natbib

\end{document}